
\subsection{Main comparison}
\label{sec:rh-main}

Table~\ref{tab:rh-main} reports \texttt{init\_params}, first
\texttt{predict} and steady-state timing for the same AlphaFold2
forward pass (\texttt{model\_3}, 118 residues, \texttt{num\_recycle=0},
float32, untrained weights) on CPU (Google Colab), GPU (NVIDIA T4,
Google Colab) and TPU (a single v5e chip on the Stanford cluster
slice). These are the original August 2026 baseline runs, each a single
call with no repetition. None of the three result
files records a timestamp or a session identifier. Archived
execution output in the August notebooks dates the CPU and GPU runs
to 2026-08-08 by the notebook's own clock; the TPU run is bounded only
by the commit that carries it. Nothing establishes that the three
backends ran within one session
(Section~\ref{sec:meth-provenance}); Section~\ref{sec:rh-variability}
reports what changed when we reran two of the three backends five
weeks later.

\begin{table}[H]
\centering\small
\caption{The August 2026 single-call baseline on the three backends:
one 118-residue inference per backend, untrained weights, float32. The
TPU row is a single v5e chip.}
\label{tab:rh-main}
\begin{tabular}{lrrr}
\toprule
Backend & \texttt{init\_params} (s) & First \texttt{predict} (s) & Steady state (s) \\
\midrule
CPU        & 41.99  & 271.98 & 212.113 \\
GPU (T4)   & 109.16 & 97.62  & 13.086  \\
TPU (v5e)  & 36.6   & 27.78  & 0.47    \\
\bottomrule
\end{tabular}
\end{table}

At steady state, the TPU is 451$\times$ faster than the CPU and
27.8$\times$ faster than the T4; the T4 itself is 16.2$\times$ faster
than the CPU. Each ratio compares exactly two measured points. We do
not chain any two of them into a claim about the TPU's advantage over
the CPU \emph{through} the GPU.
Figure~\ref{fig:rh-speedup} shows the same three timings on a log
scale, with each backend's speedup over the CPU; the TPU row carries
the single-chip caveat of Section~\ref{sec:rh-onechip}.

\subsection{TPU figures are single-chip figures}
\label{sec:rh-onechip}

The Stanford cluster slice exposes 8 TPU chips, but the single-query
path used for every number in Table~\ref{tab:rh-main} places work on
only one of them: the other seven report 0\,MB of device memory in
use throughout the run. The TPU numbers above therefore compare one
v5e chip against one T4 GPU or a two-vCPU Colab runtime (the vCPU
count is from our write-up, not a recorded field), not an 8-chip slice against
a single device. This caveat applies wherever these ratios are
cited later in the paper, including the central claim in
Section~\ref{sec:introduction}. Section~\ref{sec:parallelism} is about
what changes once the other seven chips are put to work.

\begin{figure}[H]
  \centering
  \includegraphics[width=0.75\linewidth]{speedup_hardware.pdf}
  \caption{Steady-state time per AlphaFold2 inference call (log scale),
  with the speedup of each backend over the CPU. The TPU value uses one of the slice's eight chips
  (Section~\ref{sec:rh-onechip}). All values are single-call August
  2026 baselines; Section~\ref{sec:rh-variability} reports how CPU and
  GPU timings changed across sessions.}
  \label{fig:rh-speedup}
\end{figure}

\subsection{The cold-start regime, corrected for profiler overhead}
\label{sec:rh-coldstart}

The first \texttt{predict} call is far more expensive than steady
state (271.98\,s vs.\ 212.113\,s for CPU, 97.62\,s vs.\ 13.086\,s for
GPU) because it includes JIT compilation of the full computation
graph (Section~\ref{sec:background}), not just execution. Read at face
value this gives cold/steady ratios of 1.28$\times$ (CPU) and
7.46$\times$ (GPU). Both first-\texttt{predict} timings, however, were
captured with \texttt{jax.profiler.trace} running inside the timed
region, and trace finalization itself measurably inflates the timer:
36.00\,s on CPU and 42.02\,s on GPU
(Section~\ref{sec:profiling}). Subtracting that overhead gives
corrected cold/steady ratios of approximately 1.11$\times$ for CPU and
4.25$\times$ for GPU: the compilation cost is real, but roughly
40\% smaller than the raw numbers suggest for GPU, and close to
negligible for CPU. The equivalent TPU correction cannot be applied:
no TPU run log recording profiler overhead separately is in the
repository, so the TPU's raw 59.1$\times$ cold/steady ratio is neither
confirmed nor corrected, and we do not report a corrected TPU figure.

\subsection{Cross-session variability, on two backends}
\label{sec:rh-variability}

Two of the three backends were rerun five weeks later, on
2026-09-15, with the same input and the same metric definitions as the
August baseline, from a committed revision of the benchmark script (the
August runs used an uncommitted copy of the script, so code identity
with August cannot be verified); the CPU backend was rerun twice more,
on 2026-09-16 and 2026-09-17. No rerun matches the August baseline
(Figure~\ref{fig:rh-variability}).

The CPU steady state rose from 212.113\,s (n=1, August) to
$348.861 \pm 4.690$\,s (n=5, 2026-09-15), $351.891 \pm 4.303$\,s
(n=5, 2026-09-16) and $359.273 \pm 2.787$\,s (n=5, 2026-09-17). The
three September sessions span 3\% and all sit 1.64--1.69$\times$ above
the August value. All three ran on the same host CPU model (an Intel
Xeon at 2.20\,GHz, family 6, model 79) with the same JAX version
(0.10.2); for the August session neither the host CPU model nor the JAX
version was recorded, so the gap cannot be attributed to either. With
four session-level data points the pattern looks less like symmetric
noise around a stable mean and more like a level shift between August
and September; four sessions are not enough to say which of the two
levels is typical for this Colab tier. The GPU steady state moved the other direction, from
13.086\,s (n=1, August) to $6.567 \pm 0.057$\,s (n=5, 2026-09-15), an
almost 2$\times$ \emph{drop}; there is only one September data point for
this backend. In neither case does the repository record a verified cause:
for GPU specifically, both sessions report identical GPU model,
driver and compute capability; the script copy and the numpy version
differ, while the August session's JAX version and host CPU were never
recorded, so those cannot be compared at all. No rerun isolates any of
these as the source of the gap. We report all
session values rather than picking one as canonical.

\begin{figure}[H]
  \centering
  \includegraphics[width=\linewidth]{session_variability.pdf}
  \caption{Steady-state time per call across Colab sessions for the
  same benchmark and input. (a) CPU: the August baseline (one call) and
  three September sessions (five calls each, mean and sample standard
  deviation). (b) GPU (T4): the August baseline and one September
  session. Neither backend reproduced its August value; the two moved in
  opposite directions.}
  \label{fig:rh-variability}
\end{figure}

Prior reports of cloud benchmark instability put our gaps in context,
though they do not fully account for them. Surveying that literature,
\citet{bulej2020duet} report differences of around 15\% among
instances sharing a CPU type, and similar differences for one instance
over time, while differences as large as 280\% appear between
\emph{different} CPU types. Those figures are examples from particular
studies, not bounds, so they do not tell us which case our CPU gap
(August vs.\ the September sessions, 64--69\%) and GPU gap
(approximately 100\%) fall into. A host reassignment between sessions,
which the provider does not expose to the tenant, would be consistent
with both gaps; so would a change in
the resolved software environment, which the August runs did not
record. Neither can be confirmed: the September sessions record their
host CPU and package versions, the August ones do not, so no matched
comparison is possible.
What is unusual about
our case is not the magnitude but that it appears on two different
backends of the same benchmark, run weeks apart. The two are not
independent observations: both are Google Colab runtimes, sharing a
provider and a tenancy model, and the evidence behind them is thin:
one August call per backend, three September CPU sessions and one
September GPU session. We read this as consistent with a phenomenon
the cloud-systems literature already documents for general compute,
not as independent confirmation of it on ML accelerators
(Section~\ref{sec:limitations}). Every ratio in this paper that is
derived from the August GPU baseline (B-10, B-11 and B-14 in our
results catalogue, and with them the GPU cost figures C-06 and C-09 of
Section~\ref{sec:cost}) inherits this open gap, and we did not recompute them against
the September value.
